Our results indicate the complementarity and effectiveness of multi-rate sampling and hierarchical interpolation for long-horizon time-series forecasting. Table~\ref{table:ablation_nhits_contributions} indicates that these components enforce a useful inductive bias compared to both the free-form model $\ours_4$ (plain fully connected architecture) and the parametric model \NBEATSi\  (polynomial trend and sinusoidal seasonality used as basis functions in two respective stacks). The latter obviously providing a detrimental inductive bias for long-horizon forecasting. Notwithstanding our current success, we believe we barely scratched the surface in the right direction and further progress is possible using advanced multi-scale processing approaches in the context of time-series forecasting, motivating further research.

\ours\ outperforms SoTA baselines while simultaneously providing an interpretable non-linear decomposition. Fig.~\ref{fig:motivation} and~\ref{fig:decomposition} showcase \ours\ perfectly specializing and reconstructing latent harmonic signals from synthetic and real data respectively. This novel \emph{interpretable} decomposition can provide insights to users, improving their confidence in high-stakes applications like healthcare. Finally, \ours\ hierarchical interpolation can be explored from the multi-resolution analysis perspective \citep{daubechies1992ten}. Replacing the sequential projections from the interpolation functions onto these Wavelet induced spaces is an interesting line of research.

Our study raises a question about the effectiveness of the existing long-horizon multi-variate forecasting approaches, as all of them are substantially outperformed by our univariate algorithm. If these approaches underperform due to problems with overfitting and model parsimony at the level of marginals, it is likely that the integration of our approach with Transformer-inspired architectures could form a promising research direction as the univariate results in Appendix F suggest. However, there is also a chance that the existing approaches underperform due to their inability to effectively integrate information from multiple variables, which clearly hints at possibly untapped research potential in this area. Whichever is the case, we believe our results provide a strong guidance signal and a valuable baseline for future research in the area of long-horizon multi-variate forecasting.